\chapter*{Introduction générale}
\addcontentsline{toc}{chapter}{Introduction générale} % to include the introduction to the table of content
\markboth{Introduction générale}{} %To redefine the section page head

La transformation numérique a rendu les organisations dépendantes de leur système d'information.
Messagerie, applications métier, accès distants, sauvegardes : la quasi-totalité des activités
repose sur une chaîne technique dont l'interruption, même brève, se traduit par un arrêt de
production. Longtemps considéré comme une activité de second plan, le support informatique est
devenu une fonction critique.

Cette évolution s'accompagne des progrès rapides du traitement automatique du langage naturel. Les
modèles de langage pré-entraînés savent aujourd'hui classer un texte court et bruité avec une
fiabilité utile en production, et les architectures de génération augmentée par récupération,
désignées par le sigle RAG, permettent d'interroger un historique documentaire en langage naturel
en restituant ses sources.

Le support informatique interne s'y prête particulièrement. Un service support traite peu de
problèmes inédits ; il traite surtout des variantes de cas déjà rencontrés, et chaque ticket laisse
une trace : un capital de connaissance, en principe. Ce capital reste largement inexploité, faute
d'un outil capable de qualifier les demandes et de retrouver les cas comparables. S'y ajoute une
contrainte de confidentialité, les tickets contenant des données nominatives et des éléments
d'architecture réseau.

C'est dans ce contexte que s'inscrit ce projet de fin d'études, réalisé au sein
d'OneTech Business Solutions autour de la plateforme HelpDesk AI. L'objectif est
de classer automatiquement les demandes entrantes, de retrouver dans l'historique les résolutions
comparables et d'assister la formulation de la réponse, sans qu'aucune donnée ne quitte
l'infrastructure de l'entreprise. Le service support concerné est mutualisé : onze consultants
traitent les demandes de huit entités du groupe, et l'historique de leurs interventions, soit
1\,800 tickets répartis en huit catégories, forme la matière première du travail.

Les travaux couvrent toute la chaîne, de l'analyse du besoin métier jusqu'à l'industrialisation de
la solution. Le principe directeur est le maintien de la décision humaine : l'intelligence
artificielle propose, le consultant tranche.

Le présent rapport est structuré en sept chapitres. Le premier présente l'organisme d'accueil et la
démarche de conduite retenue. Le deuxième réunit les concepts mobilisés et la spécification des
besoins. Le troisième décrit la constitution et la gouvernance du jeu de données. Le quatrième
compare plusieurs approches de classification. Le cinquième construit le système d'assistance à la
résolution. Le sixième porte sur le développement de l'application. Enfin, le septième expose
l'industrialisation et la validation de la solution.
